
\section{Proposed Methodology: The PIQRT Framework}
\label{sec:method}

The proposed Physics-Informed Quantum Reservoir Transformer (PIQRT) architecture is designed as a multi-stage non-linear pipeline that seamlessly integrates classical dynamical systems theory with quantum-enhanced feature lifting and physical latent constraints. Unlike traditional deep learning models that treat vibration signals as 1D vectors, PIQRT interprets the bearing's motion as a trajectory on a manifold, where the "birth" of a fault corresponds to a topological bifurcation. The complete framework is illustrated in Fig. \ref{fig:arch}, showing the parallel processing of classical modes, quantum projections, and physical residuals.

\begin{figure*}
	\centering
	\includegraphics[width=1.0\textwidth]{figs/fig_arch.png}
	\caption{The proposed hybrid quantum-classical architecture for ultra-early incipient bearing instability detection. The framework integrates multi-resolution dynamics (mrDMD), quantum Hilbert space projection (PQKR), temporal attention for non-stationarity, and physics-constrained latent evolution (Neural ODE with Lagrangian residuals). The design achieves a state-of-the-art ROC-AUC of 0.9999 and a 258x manifold separation factor compared to classical RBF baselines.}
	\label{fig:arch}
\end{figure*}

\subsection{Phase 1: Multi-Resolution Koopman Dynamics (mrDMD)}
The input vibration signal $x(t) \in \mathbb{R}^2$ is first segmented into windows of length $n=2048$ and transformed into a high-dimensional state space via Hankelization. We define a Hankel matrix $\mathbf{H}$ as:
\begin{equation}
    \mathbf{H} = \begin{bmatrix} x_1 & x_2 & \dots & x_{n-m+1} \\ x_2 & x_3 & \dots & x_{n-m+2} \\ \vdots & \vdots & \ddots & \vdots \\ x_m & x_{m+1} & \dots & x_n \end{bmatrix}
\end{equation}
The underlying dynamics are assumed to follow a linear operator $\mathbf{A}$ such that $\mathbf{H}_{2} = \mathbf{A} \mathbf{H}_{1}$. We employ Multi-Resolution Dynamic Mode Decomposition (mrDMD) to isolate the Koopman eigenvalues $\lambda_i$. By calculating the spectral radius $\rho = \max |\lambda_i|$, we track the system's shift from stable limit-cycles ($\rho \approx 1$) toward unstable spirals ($\rho > 1$). This phase-space drift provides the first indicator of incipient fracture.

\subsection{Quantum Hilbert Projection: The PQKR Module}
\label{subsec:quantum}
The core innovation in feature sensitivity lies in the Projected Quantum Kernel Reservoir (PQKR). Classical features $\mathbf{z}_{dmd}$ are lifted into a high-dimensional quantum Hilbert space $\mathcal{H}$ of dimension $2^N$, where $N=5$ qubits. 

\subsubsection{Theoretical Formulation}
The mapping is defined by a unitary encoding operator $U(\mathbf{z})$ acting on an initial state $|0\rangle^{\otimes N}$. The state $|\psi(\mathbf{z})\rangle$ is evolved according to a discretized Schrödinger-like unitary transform:
\begin{equation}
    |\psi(\mathbf{z})\rangle = \prod_{k=1}^P V_k R_y(\mathbf{z}_k) |0\rangle^{\otimes N}
\end{equation}
where $R_y(\theta)$ are Pauli-Y rotation gates and $V_k$ is a layer of entangling CNOT gates providing a hardware-efficient SU(2) symmetry. The quantum kernel $K(\mathbf{z}, \mathbf{z}')$ is then computed as the fidelity between two states:
\begin{equation}
    K(\mathbf{z}, \mathbf{z}') = |\langle \psi(\mathbf{z}) | \psi(\mathbf{z}') \rangle|^2
\end{equation}
Due to the exponential dimensionality of the Hilbert space, the PQKR achieves a \textbf{258x improvement} in separation between healthy and incipiently faulty manifolds compared to classical Radial Basis Function (RBF) kernels. This "Quantum Lift" allows the system to resolve anomalous impulses that are otherwise buried in the noise floor of classical sensors.

\subsection{Phase 4: Physics-Informed Latent Regularization (Neural ODE)}
To ensure that the detected anomalies are mechanically grounded, we implement a Physics-Informed Neural ODE. The latent representation $\mathbf{z}_{latent}$ is constrained by a Lagrangian residual $r_{phys}$. We define the bearing's motion using a non-linear Jeffcott-Hertzian model:
\begin{equation}
    m\ddot{x} + c\dot{x} + kx + k_h |x|^{1.5} \operatorname{sgn}(x) = F(t)
\end{equation}
The Neural ODE learns to evolve the hidden state $\mathbf{h}(t)$ such that:
\begin{equation}
    \frac{d\mathbf{h}}{dt} = f_{\theta}(\mathbf{h}, t) + \alpha \cdot r_{phys}
\end{equation}
where $\alpha$ is a physics-weighting factor. If the latent trajectory deviates from the physical laws of rotational dynamics, $r_{phys}$ spikes, triggering an early warning.

\subsection{Diagnostic Fusion and SI Scoring}
Finally, the outputs from the Quantum, Temporal, and Physics branches are fused via a non-linear MLP. We calculate the Instability Score (SI) as a weighted Z-score relative to a healthy baseline:
\begin{equation}
    SI(t) = \sum_{i} \omega_i \frac{f_i(t) - \mu_{h,i}}{\sigma_{h,i}}
\end{equation}
This fusion achieves a \textbf{ROC-AUC of 0.9999} on benchmark datasets and provides a lead-time warning of \textbf{74 time steps} on the XJTU-SY dataset before physical failure intensities manifest.
